\documentclass[10pt,a4paper]{amsart}
\usepackage[margin=19mm]{geometry}
\usepackage{amsmath,amssymb,mathtools}
\usepackage[scr=rsfs,cal=euler]{mathalfa}
\usepackage{xfrac}
\usepackage[unicode,hidelinks,bookmarks=false]{hyperref}
\newcommand{\ie}{i.e.\ }
\newcommand{\cf}{cf.\ }
\newcommand{\resp}{respectively\ }
\newcommand{\Subsection}{\subsection}
\providecommand{\nobreakdash}{\nobreak\hbox{-}}
\setlength{\emergencystretch}{2em}
\multlinegap0pt
\usepackage{bm}
\usepackage{bbm}
\usepackage{dsfont}
\usepackage{xspace}
\usepackage{cancel}
\usepackage{mathtools}
\usepackage{amsthm}
\makeatletter
\@ifundefined{theorem}{\newtheorem{theorem}{Theorem}}{}
\@ifundefined{lemma}{\newtheorem{lemma}[theorem]{Lemma}}{}
\makeatother
\makeatletter
\def\a{\alpha}
\def\eps{\varepsilon} 
\expandafter\ifx\csname proof\endcsname\relax
\newenvironment{proof}{\noindent{\bf Proof:}}{\qed}
\fi
\makeatother
\title[On likelihood of a Condorcet winner for uniformly random and]{On likelihood of a Condorcet winner for uniformly random and independent voter preferences}
\author{Boris Pittel}
\date{Source: arXiv:2506.20613v5 (CC BY 4.0)}
\begin{document}
\maketitle

\noindent\emph{Source and licence.} On likelihood of a Condorcet winner for uniformly random and independent voter preferences. Version arXiv:2506.20613v5 (CC BY 4.0), distributed under the Creative Commons Attribution 4.0 International licence, as recorded in the arXiv licence metadata for this exact version and shown on the abstract page https://arxiv.org/abs/2506.20613v5. Full rights notice: licenses/arxiv-2506.20613-CC-BY-4.0.txt.\par\smallskip

\noindent\textit{Editorial scope.} Counted unit: the Lemma whose source label is \texttt{nlem3} in the article (source0.tex lines 392--393, source label \texttt{nlem3}), delivered together with the article's own complete original proof of it (source0.tex lines 394--473, one \texttt{proof} environment of 80 lines). No mathematics is written, repaired or re-derived: statement and proof are the article's own text line for line. Required context retained uncounted: nlem2 (lines 267--273). Every definition, display and label of those lines is retained unchanged. Omitted from this package and not redistributed: 1--266, 274--391, 474--820. Nothing inside a retained excerpt is omitted or altered: every retained body is delivered exactly as the article's lines are written, so no omission receipt of any kind accompanies this package. Citations: the retained excerpts carry no citation command at all, so no bibliography entry of the article is transcribed and no reference is redistributed. Reference adaptation: none was needed and none was made; every reference inside the delivered unit resolves inside it, so no unresolved reference or citation is produced. Printed slips kept literal: no printed word, symbol, hypothesis or number of the counted statement or of its proof was altered. Layout and class: the article's own class is \texttt{amsart} with packages including xcolor, graphicx, latexsym, amsfonts, amsmath, amssymb, color; the wrapper is the lane's pinned \texttt{amsart} editorial wrapper at 10pt on A4, which restates the theorem-like environments the retained text needs and adds the article's own macro definitions that the retained text needs (\texttt{\textbackslash a}, \texttt{\textbackslash eps}), with extra wrapper packages (bm, bbm, dsfont, xspace, cancel, mathtools). The delivered numbering is the unit's own. Assets: no retained span contains a figure environment, a graphics-inclusion command, a PDF-page inclusion, a table or a TikZ picture; no image file is copied into this package, the package declares no asset dependency and no image file is redistributed with it. Licence: version arXiv:2506.20613v5 of this article is distributed under the Creative Commons Attribution 4.0 International licence, as recorded in the arXiv licence metadata for this exact version and shown on the abstract page https://arxiv.org/abs/2506.20613v5. A full rights notice is delivered at licenses/arxiv-2506.20613-CC-BY-4.0.txt. Characters: every retained excerpt is pure ASCII, so the delivered engine, XeLaTeX with its default Latin Modern text font, reports no missing character.
\begin{lemma}\label{nlem2} For $n\ll k^{-1/2}$,
\begin{equation*}
\Bbb E\bigl[G^{n-1}(y(\bold X))\bigr]=\!\int_{-\infty}^{\infty}G^{n-1}(y)\frac{e^{-y^2}}{\sqrt{\pi}}\,dy +O(nk^{-1/2}).
%\frac{\sqrt{2}}{n}\int_0^{\infty} ye^{-y^2/2}\bigl[G^{n}(-y)-G^{n}(y)\bigr]\,dy +
%O(n k^{-1/2}).
\end{equation*}
\end{lemma}

\begin{lemma}\label{nlem3} $I_n=O\bigl(\exp(-n^{\eps(n)})\bigr)$ for every $\eps(n)\downarrow 0$. In particular, $I_n$ is super-polynomially small.
\end{lemma}

\begin{proof} 
First,  
\begin{multline*}
%\int_{-\infty}^{\infty}G^{\nu}(y)\frac{e^{-y^2}}{\sqrt{\pi}}\,dy
I_{n}=\frac{\sqrt{2}}{n}
%\int_{-\infty}^{\infty}n G^{n-1}(y)\frac{e^{-y^2}/2}{\sqrt{2\pi}}\cdot e^{-y^2/2}\,dy\\
\int_{-\infty}^{\infty}(-G^n(y))' e^{-y^2/2}\,dy\\
=\frac{\sqrt{2}}{n}\int_{-\infty}^{\infty}(-y)e^{-y^2/2}G^{n}(y)\,dy\\
=\frac{\sqrt{2}}{n}\int_0^{\infty} ye^{-y^2/2}\bigl[G^{n}(-y)-G^{n}(y)\bigr]\,dy.
\end{multline*}
%First of all, $I_{\nu}=o(\nu^{-1})$, since $\int_0^{\infty}ye^{-y^2/2}\,dy=1$, and $G^{\nu+1}(-y)\to 0$ for every fixed $y$. More precisely,
Since $G(y)\le 1/2$ for $y\ge 0$, we have $
\int_0^{\infty} ye^{-y^2/2}G^{n}(y)\,dy=O(2^{-n})$. So $I_n=\frac{\sqrt{2}}{n}(J_n+O(2^{-n}))$, where 
\[
J_n:=\int_0^{\infty}ye^{-y^2/2} G^{n}(-y)\,dy=\int_0^{\infty}ye^{-y^2/2}\left(1-\frac{1}{\sqrt{2\pi}}\int_y^{\infty}e^{-z^2/2}\,dz\right)^n\,dy.
\]
Because of how the integrand's depends on $n$, we expect that the dominant contribution to $J_n$ comes from a relatively small neighborhood of the integrand's maximum point. It is easy to see that the integrand is strictly log-concave, and since the integrand approaches zero both at $y=0$ and $y=\infty$, it attains its unique maximum at a unique point. A closer look shows that the maximum
is very near $\bar y:=\left(2\log\left(\frac{n}{\sqrt{\log n}}\right)\right)^{1/2}$. We do not have to prove this clam, since its sole purpose is to 
{\it motivate\/} a rigorous argument. 
Let $J_{n,1}$ and $J_{n,2}$ denote the contributions to $J_n$ coming respectively from $y\in [0,y(\a)]$ and $y\in [y(\a),\infty)$, $y(\a):=\a \bar y$. Here $\a=\a(n)\uparrow 1$ will be chosen based on bounds on $J_{n,j}$.

{\bf (a)\/} Using a classic inequality
 \begin{equation}\label{verynew1}
\int_y^{\infty}e^{-z^2/2}\,dz\ge \bigl(y^{-1}-y^{-3}\big)e^{-y^2/2},
\end{equation}
we have: 
\begin{multline}\label{new4}
J_{n,1}\le\left(1-\frac{1}{\sqrt{2\pi}}\int_{y(\a)}^{\infty}e^{-z^2/2}\,dz\right)^n\int_0^{y(\a)}ye^{-y^2/2}\,dy\\
\le \exp\left(-\frac{n}{\sqrt{2\pi}}\int_{y(\a)}^{\infty}e^{-z^2/2}\,dz\right)\\
\le\exp\left(-\bigl(y^{-1}(\a)-y^{-3}(\a)\bigr)\frac{n}{\sqrt{2\pi}}e^{-y^2(\a)/2}\right).
\end{multline}
%\int_2^{y(\eps)} y\exp\left(-\frac{y^2}{2}-\frac{n}{\sqrt{2\pi}}\left(\frac{1}{y}-\frac{1}{y^3}\right)e^{-y^2/2}\right)\,dy\\
%\le \int_0^{y(\eps)}\!\! y\exp\left(\!-\frac{y^2}{2}-\frac{n}{\sqrt{2\pi}}\frac{3}{4y(\eps)}e^{-y^2/2}\!\right)\,dy\\
%=\int_0^{y^2(\eps)/2}\!\!\!\exp\left(\!-z-\frac{n}{\sqrt{2\pi}}\frac{3}{4y(\eps)}e^{-z}\!\right) dz\\
%=\int_{e^{-y^2(\eps)/2}}^1\exp\left(-\frac{n}{\sqrt{2\pi}}\frac{3}{4y(\eps)}\eta\right)\,d\eta\\
%\le \frac{\exp\left(-\frac{n}{\sqrt{2\pi}}\frac{3}{4y(\eps)}e^{-y^2(\eps)/2}\right)}{\frac{n}{\sqrt{2\pi}}\frac{3}{4y(\eps)}}.
%\end{multline*}
%And obviously the contribution to $J_{n,1}$ from $[0,2]$ is at most $q^n$, where $q:=1-\frac{1}{\sqrt{2\pi}}\int_2^{\infty}e^{-z^2/2}\,dz$.

{\bf (b)\/} Next,
\begin{multline*}
J_{n,2}=\int_{y(\a)}^{\infty}ye^{-y^2/2}\left(1-\frac{1}{\sqrt{2\pi}}\int_y^{\infty}e^{-z^2/2}\,dz\right)^n\,dy\\
\le \int_{y(\a)}^{\infty}ye^{-y^2/2}
\exp\left(-\frac{n}{\sqrt{2\pi}}\int_y^{\infty}e^{-z^2/2}\,dz\right)\,dy\\
=\int_{y^2(\a)/2}^{\infty}\exp\left(-u-\frac{n}{\sqrt{2\pi}}\int_{\sqrt{2u}}^{\infty}e^{-z^2/2}\,dz\right)\,du\\
=\int_{0}^{e^{-y^2(\a)/2}}e^{-nK(v)}\,dv, \quad K(v):=(2\pi)^{-1/2}\int_{\sqrt{2\log(1/v)}}^{\infty}e^{-z^2/2}\,dz.
\end{multline*}
Now, $K'(v)=\frac{1}{2\sqrt{\pi}}\log^{-1/2}(1/v)(>0)$ increases with $v$, so that $K(v)$ is convex; hence
$
K(v)\ge K(e^{-y^2(\a)/2})+K'(e^{-y^2(a)/2})\bigl(v-e^{-y^2(\a)/2}\bigr).
$
Therefore 
$
J_{n,2}\le \frac{e^{-nK(e^{-y^2(\a)/2})}}{nK'(e^{-y^2(\a)/2})}.
$
Here
\begin{align*}
nK(e^{-y^2(\a)/2})&=\frac{n}{\sqrt{2\pi}}\int_{y(\a)}^{\infty}e^{-z^2/2}\,dz\\
&\ge \bigl(y^{-1}(\a) -y^{-3}(\a)\bigr)\frac{n}{\sqrt{2\pi}}e^{-y^2(\a)/2},\\
%\ge c_2\frac{n}{\sqrt{\log n}}\cdot\exp\Bigl(-(1-\eps)^2\log\bigl(\frac{n}{\sqrt{\log n}}\bigr)\Bigr)=c_2\Bigl(\frac{n}{\sqrt{\log n}}\Bigr)^{2\eps -\eps^2}\to\infty,\\
nK'(e^{-y^2(\eps)/2})&=\frac{n}{\sqrt{2\pi}}y^{-1}(a).
\end{align*}
So, $J_{n,2}$ is at most of order
\begin{equation}\label{new5}
\frac{\exp\left(-\bigl(y^{-1}(\a) -y^{-3}(\a)\bigr)\frac{n}{\sqrt{2\pi}}e^{-y^2(\a)/2}\right)}{ny^{-1}(\a)}
\end{equation}
Here $ny^{-1}(\a)$ is exactly of order $n\log^{-1/2} n\to \infty$. So \eqref{new4} and \eqref{new5} together imply that 
\begin{multline*}
J_n=O\left[\exp\left(-\bigl(y^{-1}(\a)-y^{-3}(\a)\bigr)\frac{n}{\sqrt{2\pi}}e^{-y^2(\a)/2}\right)\right]\\
=O\left[\exp\left(-\frac{1+o(1))}{2\sqrt{\pi}}\frac{n}{\sqrt{\log n}}\left(\frac{\sqrt{\log n}}{n}\right)^{\a^2}\right)\right]\\
=O\left[\exp\left(-\frac{1+o(1))}{2\sqrt{\pi}}\left(\frac{n}{\sqrt{\log n}}\right)^{1-\a^2}\right)\right]
\end{multline*}
%\le \exp\lbigl(-\exp(\log^{-s}n)\bigr),
For $\a:=1- \eps(n)$, $\eps(n)\downarrow 0$, we have $1-\a^2=2\eps(n)-\eps^2(n)$, whence $J_n\le \exp\bigl(-n^{\eps(n)}\bigr)$.
We conclude that
\[
I_n=\frac{\sqrt{2}}{n}\bigl(J_n+O(2^{-n})\bigr)=O\bigl[\exp\bigl(-n^{\eps(n)})\bigr].
\]
In particular, selecting $\eps(n)=\frac{\log(\log(\ell\log n))}{\log n}$, $\ell>0$ being fixed. we obtain $I_n=O(n^{-\ell})$. This completes the proof of Lemma \ref{nlem2}.
\end{proof}

\end{document}
